\centerline{\bf \Huge Оценка + пример I}

\begin{flushright}
        \scriptsize{Никогда не думайте, что вы уже всё знаете. И как бы высоко ни оценивали вас, всегда имейте мужество сказать себе: «Я невежда».}
\end{flushright}

\problem{\textbf{1}} Какое наибольшее количество ладей можно поставить на шахматную доску так, чтобы они не били друг друга? (ладья бьет по вертикали и горизонтали)

\begin{flushright} \textbf{2 балла}\\ 
\end{flushright}

\problem{\textbf{2}} Найдите наименьшее натуральное число кратное 5, сумма цифр которого равна 25.

\problem{\textbf{3}} На клетчатой доске 100×100 закрасили Х прямоугольников 1×2.
Оказалось, что в каждой строке и в каждом столбце есть хотя
бы одна закрашенная клетка. При каком наименьшем Х это
возможно?

\begin{flushright} \textbf{4 балла}\\ 
\end{flushright}

\problem{\textbf{4}} Найдите наибольшее число, любые две последовательные цифры которого образуют точный квадрат в том порядке, в котором они стоят. (Точный квадраты это: 0, 1, 4, 9, 16, 25...)

\begin{flushright} \textbf{6 баллов}\\ 
\end{flushright}

\problem{\textbf{5}} Пятизначное число называется неразложимым, если оно не
раскладывается в произведение двух трёхзначных чисел. Какое
наибольшее количество неразложимых пятизначных чисел может идти
подряд?

\begin{flushright} \textbf{8 баллов}\\ 
\end{flushright}
